\documentclass[12pt,a4paper,oneside]{article}
\usepackage[brazil]{babel}
\usepackage[utf8]{inputenc}
\usepackage{olymp}

\setcounter{problem}{0}

\begin{document}

\begin{problem}{Quadrado mágico}{quadrado}{2}

Senhor Coelho é conhecido mundialmente pela fabricação de quadrados mágicos de dimensões $3 \times 3$. Um quadrado é chamado mágico quando a soma dos elementos de uma determinada linha, coluna ou diagonal é sempre igual.

Infelizmente, assaltantes invadiram recentemente a oficina do Sr. Coelho e roubaram alguns dos números de seus quadrados mágicos. Felizmente os meliantes não conseguiram roubar mais do que 3 números de cada quadrado. Desesperado, pois devia entregar os quadrados naquele dia, o Sr. Coelho veio procurar a sua ajuda para tentar completar os quadrados com os números faltantes.

Escreva um programa que, dado um quadrado mágico com alguns números faltando, determine qual era o quadrado mágico original.


%\Notes
%
%Observações, se tiver.

\InputFile

A entrada contém três linhas, cada uma contendo três inteiros entre 0 e 2 x ${10^4}$, inclusive. O número zero representa os dígitos que foram roubados. Existem no máximo três números zero na entrada.

\OutputFile

Seu programa deve imprimir três linhas, cada uma contendo três inteiros, descrevendo a configuração original do quadrado mágico.

% \Notes
% É garantido que sempre existe uma divisão que satisfaz as condições dos reinos.

\Example

\begin{example}
\exmp{
4 9 2
3 0 7
8 1 6
}{
4 9 2
3 5 7
8 1 6
}%
\end{example}

\begin{example}
\exmp{
0 12 12
16 10 0
8 8 14
}{
6 12 12
16 10 4
8 8 14
}%
\end{example}

\begin{example}
\exmp{
495 468 0
0 522 414
441 0 549
}{
495 468 603
630 522 414
441 576 549
}%
\end{example}

%Comentários sobre o exemplo, se necessário.

\end{problem}

\end{document}
